\vspace{2.5pt plus 3pt minus 0.5pt}\noindent\hspace*{1.2em}\kbvnum{(೯)}\hspace{0.55em}ಬೆಳಕು ಬರಲು ಕತ್ತಲೆ ಹೋಯಿತು.\par

\vspace{2.5pt plus 3pt minus 0.5pt}\noindent\hspace*{1.2em}\kbvnum{(೧೦)}\hspace{0.55em}ಅವಳು ಬೆಂಕಿಯ ಜ್ವಾಲೆಯಲ್ಲಿ ಧಗಧಗ ಉರಿದು ಹೋದಳು.\par

\vspace{2.5pt plus 3pt minus 0.5pt}\noindent\hspace*{1.2em}\kbvnum{(೧೧)}\hspace{0.55em}ಕಾದ ಎಣ್ಣೆಯಲ್ಲಿ ಸಾಸಿವೆ ಚಟಚಟ ಎನ್ನುವುದು.\par

\vspace{2.5pt plus 3pt minus 0.5pt}\noindent\hspace*{1.2em}\kbvnum{(೧೨)}\hspace{0.55em}ರಾಮನೂ ಸೀತೆಯೂ ಲಕ್ಷ್ಮಣನೂ ಕಾಡಿಗೆ ಹೋದರು.\par

\vspace{2.5pt plus 3pt minus 0.5pt}\noindent\hspace*{1.2em}\kbvnum{(5)}\hspace{0.55em}ಎರಡು ವಾಕ್ಯಗಳನ್ನು ಜೋಡಿಸಲು “ಆದ್ದರಿಂದ” ಎಂಬ ಅವ್ಯಯವನ್ನು ಹೇಗೆ ಬಳಸಬಹುದೆಂಬುದಕ್ಕೆ ಒಂದು ಉದಾಹರಣೆ ಕೊಡಿ.\par

\vspace{2.5pt plus 3pt minus 0.5pt}\noindent\hspace*{1.2em}\kbvnum{(6)}\hspace{0.55em}ಅವಧಾರಣಾರ್ಥಕಾವ್ಯಯ ಎಂದರೇನು? ಎರಡು ಉದಾಹರಣೆ ಗಳನ್ನು ಕೊಡಿ.\par

\kbvchapter{ಪಾಠ ೧೭}{ದ್ವಿರುಕ್ತಿ}

\indent ದ್ವಿರುಕ್ತಿ ಎಂದರೆ ಎರಡು ಬಾರಿ ಹೇಳುವುದು ಎಂದರ್ಥ. ಸಂದರ್ಭಾನು ಸಾರವಾಗಿ ಸನ್ನಿವೇಶಕ್ಕೆ ಒತ್ತುಕೊಡುವ ಉದ್ದೇಶದಿಂದ ಒಂದೇ ಶಬ್ದವನ್ನು ಎಡಬಿಡದೆ ಎರಡು ಬಾರಿ ಹೇಳುವುದು ರೂಢಿಯಲ್ಲಿದೆ. ಇದನ್ನು ದ್ವಿರುಕ್ತಿ ಎನ್ನುವರು.\par

\vspace{3pt plus 2pt}\noindent\hspace*{1.2em}\kbvlabel{ಉದಾ}\par\vspace{1pt}

\noindent\hspace*{2.2em}ರೈಲು ಹತ್ತಲು ಅವನು ಓಡಿಓಡಿ ಬಂದನು.\par

\noindent\hspace*{2.2em}ಹೌದು ! ಹೌದು ! ಅವನೇ ಸರ ಕದ್ದವನು.\par

\noindent\hspace*{2.2em}ಬೇಡಬೇಡ, ನನಗೆ ನಿನ್ನ ಸಹವಾಸವೇ ಬೇಡ.\par

\noindent\hspace*{2.2em}ತಿಂದೂತಿಂದೂ ಸಾಕಾಯಿತು.\par

\noindent\hspace*{2.2em}ನೀನು ನನ್ನನ್ನು ಆಗಾಗ ಬಂದು ಕಾಣುತ್ತಿರು.\par

\noindent\hspace*{2.2em}ಊರೂರು ಸುತ್ತಿ ನನಗೆ ಆಯಾಸವಾಗಿದೆ.\par

\indent ಅಬ್ಬಬ್ಬಾ ! ಎಷ್ಟು ದೊಡ್ಡ ಹಾವು. ಇಲ್ಲಿ ಬನ್ನಿ ಇಲ್ಲಿ ಬನ್ನಿ, ನನ್ನ ಮಾತು ಕೇಳಿ. ಈ ಮಕ್ಕಳಿಗೆ ವಿವೇಕ ಹೇಳಿಹೇಳಿ ಸಾಕಾಯಿತು.\par

\noindent\hspace*{2.2em}ಅತ್ತೂ, ಅತ್ತೂ ಕಣ್ಣು ಕೆಂಪಗಾಯಿತು.\par

\indent ಈ ವಾಕ್ಯಗಳಲ್ಲಿ ದಪ್ಪಕ್ಷರಗಳಲ್ಲಿರುವ ಪದಗಳು ದ್ವಿರುಕ್ತಿಗಳು.\par
